% concatenated from CSDunitary.tex
\documentclass[a4paper,12pt]{article}
\usepackage[left=3.17cm,right=3.17cm,top=2.54cm,
headheight=0.5cm,headsep=0.54cm,bottom=2.54cm,footskip=0.79cm
]{geometry}

\usepackage{amsmath,amsthm,amssymb,cite}

\usepackage{hyperref}


\newtheorem{corollary}{Corollary}
\newtheorem{definition}{Definition}
\newtheorem{lemma}{Lemma}
\newtheorem{theorem}{Theorem}



\begin{document}

\title{A note on unitary invariance of Connes spectral distances of quantum states}

\author{Bing-Sheng Lin$^{1,2,}$\thanks{Corresponding author. 
		E-mail address: sclbs@scut.edu.cn}, Ji-Hong Wang$^{3}$, Zhi-Kang You$^{1}$\\
	\small $^{1}$ School of Mathematics, South China University of Technology,
		Guangzhou 510641, China\\
	\small $^{2}$ Institute for Mathematics, Astrophysics and Particle Physics,\\
	\small  Radboud University Nijmegen, Heyendaalseweg 135, 6525 AJ Nijmegen, The Netherlands\\
	\small $^{3}$ Guangzhou University of Software, Guangzhou 510990, China
}


\date{\today}


\maketitle

\begin{abstract}
In this paper, we study the properties of Connes spectral distances between quantum states under unitary transformations. We mainly focus on spectral triples with matrix algebras acting on finite dimensional Hilbert spaces. We prove that there are some finite spectral triples in which the Lipschitz seminorms are equal to the operator norms. We also explicitly construct some spectral triples in which the Connes spectral distances between quantum states are exactly the quantum trace distances.
These results are helpful for us to better study the relationships among Connes spectral distance, quantum trace distance and other quantum distance measures. These concrete examples are significant for studies of geometric structures of finite spectral triples and mathematical relations of qubits and other quantum states in the framework of noncommutative geometry.\\
  
\textit{Keywords}: Connes spectral distance; Unitary invariance; Quantum trace distance.
  
%\textit{PACS}: 02.40.Gh, 03.65.Fd, 03.65.Aa

\end{abstract}



\section{Introduction}
In the 1980’s, Connes formulated the mathematically rigorous framework of noncommutative geometry \cite{Connes}.
The main objects in noncommutative geometry are spectral triples, which can be regarded as noncommutative versions of normal spaces.
Due to the noncommutativity, there is no traditional point in a spectral triple. This is different from normal commutative spaces.
In a spectral triple, a pure state is the
analog of a traditional point in a normal commutative space.
So there is no normal concept of distance between two points in a noncommutative space, but
we can calculate some kinds of distance measure between the states, such as the
Connes distance \cite{Connes1}.
The Connes spectral distance between pure states corresponds to the geodesic distance
between points.
The Connes spectral distances in some kinds
of noncommutative spaces have already been studied in literature
\cite{Bimonte,Dai,Cagnache,Wallet,Martinetti1,Martinetti,DAndrea,Pythagoras,Pythagoras1,Franco,Scholtz,Chaoba,Revisiting,Kumar,Barrett,Chakraborty,Lin1,Lin2,Clare}.
For example, Cagnache \textit{et. al.} computed
Connes spectral distances between the pure states which corresponds to
eigenfunctions of the quantum harmonic oscillators in the Moyal
plane\cite{Cagnache}. Martinetti \textit{et. al.} obtained the
spectral distance between coherent states in the so-called double Moyal plane
\cite{Martinetti}. D’Andrea \textit{et. al.} have studied the Pythagoras’
theorem in noncommutative geometry \cite{Pythagoras}. They showed that for non-pure states
it is replaced by some Pythagoras inequalities. Franco and Wallet also studied
metrics and causality on Moyal planes \cite{Franco}. Scholtz and his collaborators have studied the Connes
spectral distances of harmonic oscillator states and also coherent states in
Moyal plane and fuzzy space \cite{Scholtz,Chaoba,Revisiting}.
Kumar \textit{et. al.} used Dirac eigen-spinor method to compute spectral distances in doubled Moyal plane \cite{Kumar}. Chakraborty \textit{et. al.} also studied the spectral
distance on Lorentzian Moyal plane \cite{Chakraborty}.
Lin \textit{et. al.} studied the relation between Connes spectral
distance and the noncommutativity of phase spaces \cite{Lin2}.

In quantum physics, most of the properties of physical systems are invariant
under unitary transformations. Connes spectral distances are also some essential
physical properties of quantum states. Therefore, it is interesting to study
some types of Connes spectral distances which are invariant under unitary transformations of the states.
This is also a special kind of isometry of the state space being studied in some other literature, such as Ref.~\cite{Martinettip}.
The unitary invariance of spectral distances can be regarded as a natural property of a spectral triple.
Quantum trace distance is one of the most useful distance measures between
quantum states in quantum information sciences \cite{Nielsen}. It is also unitary invariant.
So it is also very significant to study spectral triples with Connes spectral
distances which equal quantum trace distances.

In the present work, we will study some properties and concrete examples of Connes spectral distances between quantum states under unitary transformations.
This paper is organized as follows. In Sec.~\ref{sec2}, we review the definitions of spectral triples and Connes spectral distances.
In Sec.~\ref{sec3}, we review some elementary properties of Connes spectral distances and optimal elements.
In Sec.~\ref{sec4}, we study the unitary invariance of Connes spectral distances
of quantum states. In Sec.~\ref{sec5}, we prove that there are some finite spectral triples in which the Lipschitz seminorm is equal to the operator norm, and then the Connes spectral distances between states are exactly the quantum trace distances.
In Sec.~\ref{sec6}, we explicitly construct some spectral triples equip with quantum trace distances between one-qubit states. 
Some discussions and conclusions are given in Sec.~\ref{sec7}.



\section{Connes spectral distances in spectral triples}\label{sec2}

In the framework of noncommutative geometry \cite{Connes}, a noncommutative version of space corresponds to a spectral triple $(\mathcal{A}, \mathcal{H},
\mathcal{D})$, where $\mathcal{A}$ is an involutive algebra acting on a
Hilbert space $\mathcal{H}$ through a representation $\pi$, and the Dirac
operator $\mathcal{D}$ is a self-adjoint, densely defined operator on
$\mathcal{H}$ which satisfies:

1.\;\;$\mathcal{D}$ can be unbounded operator but $[\mathcal{D},
\pi (a)]$ is bounded with $a \in \mathcal{A}$;

2.\;\;$\mathcal{D}$ has compact resolvent, for $\lambda \in
\mathbb{C} \backslash \mathbb{R}$, $(\mathcal{D}- \lambda)^{- 1}$ is compact
when the algebra $\mathcal{A}$ is unital or $\pi (a)  (\mathcal{D}-
\lambda)^{- 1}$ be compact if it is non-unital.

In the present work, we will consider finite matrix
algebras $\mathcal{A}=\mathbb{M}_n (\mathbb{C})$ with representations $\pi$ on
finite dimensional Hilbert spaces $\mathcal{H}= \bigoplus_{i = 1}^m
\mathbb{C}^n$, $n \geqslant 2$, $m \geqslant 1$, and $\pi (a) \in
\mathbb{M}_{mn} (\mathbb{C}) = L (\mathcal{H})$.
In most of the cases, we always assume that the
representation $\pi$ is unital, namely, $\pi (\mathbb{I}_{n}) =\mathbb{I}_{mn}$. For example, $\pi (a) =\mathbb{I}_m \otimes
a$. So for any unitary element $U \in \mathcal{A}$, $\pi (U)$ is also unitary.

A state $\omega$ on the algebra $\mathcal{A}$ is a linear functional $\omega :
\mathcal{A} \rightarrow \mathbb{C}$ which is positive i.e. $\omega (a^{\ast}
a) \geq 0, \forall a \in \mathcal{A}$ and has a norm $1$. For the quantum
state $\omega$ which is normal and bounded, it can also be represented by the
density operator $\rho \in \mathcal{A}$, which is positive semi-definite,
self-adjoint and trace $1$. The action of the state $\omega$ on an element $e
\in \mathcal{A}$ can be written as \cite{Revisiting}
\begin{equation}
  \omega (e) = \mathrm{tr} (\rho e) .
\end{equation}
So the state $\omega$ can be represented by the corresponding density matrix
$\rho$. For convenience, we will denote the state $\omega$ by the
corresponding density matrix $\rho$ in the following content if there is no
confusion.

Suppose the quantum states $\omega_1$ and $\omega_2$ correspond to the density
matrices $\rho_1$ and $\rho_2$, respectively.

\begin{definition}
  \cite{Connes1} The Connes distance between quantum states $\omega_1$ and
  $\omega_2$ is
  \begin{equation}
    d (\omega_1, \omega_2) \equiv d (\rho_1, \rho_2) = \sup_{e \in B} |
    \mathrm{tr} (\rho_1 e) - \mathrm{tr} (\rho_2 e) | = \sup_{e \in B} |
    \mathrm{tr} (\Delta \rho \, e) |,
  \end{equation}
  where $\Delta \rho = \rho_1 - \rho_2$, and the set
  \begin{equation}
    B := \{e \in \mathcal{A}: \|[\mathcal{D}, \pi (e)]\|_{op} \leqslant 1\},
  \end{equation}
  where $\|A\|_{op}$ is the operator norm of $A$,
  \begin{equation}
    \|A\|_{op} \equiv \sup_{\psi \in \mathcal{H}, \| \psi \| = 1} \|A \psi \|,
    \qquad \|A\|^2 \equiv \mathrm{tr} (A^{\dag} A) = \langle A, A \rangle,
  \end{equation}
  and $\langle \cdot, \cdot \rangle \equiv \langle \cdot, \cdot
  \rangle_{\mathrm{HS}}$ is the Hilbert-Schmidt inner product, $\langle A, B
  \rangle \equiv \langle A, B \rangle_{\mathrm{HS}} := \mathrm{tr} (A^{\dag}
  B)$.
\end{definition}

The inequality $\|[\mathcal{D}, \pi (e)]\|_{op} \leqslant 1$ is the so-called
\textit{ball condition}. The map $L (a) = \|[D, \pi (a)]\|_{op}$ is
a Lipschitz seminorm.

Obviously, if $e \in B$, then we also have $- e \in B$, so we also have
\begin{equation}
  d (\rho_1, \rho_2) = \sup_{e \in B} [\mathrm{tr} (\rho_1 e) - \mathrm{tr}
  (\rho_2 e)] = \sup_{e \in B} [\mathrm{tr} (\Delta \rho \, e)] =
  \mathrm{tr} (\Delta \rho \, e_o),
\end{equation}
where $e_o$ is the so-called \textit{optimal element}.
Note that the last equality above makes sense only if the supremum is attained (for simplicity, we only consider these cases in the present work).
Since
Hermitian elements can give the supremum in the Connes spectral distance
functions \cite{Iochum}, one only need to consider the optimal elements
$e_o$ being Hermitian.

Regarding that $(\Delta \rho)^{\dag} = \Delta \rho$, we also have
\begin{equation}
  d (\rho_1, \rho_2) = \sup_{e \in B} (\mathrm{tr} [(\Delta \rho)^{\dag} 
  \, e]) = \sup_{e \in B} \langle \Delta \rho, e \rangle =
  \langle \Delta \rho, e_o \rangle .
\end{equation}
Obviously, there is $d (\rho_1, \rho_2) \geqslant 0$. In the present work, we
only consider the cases $d (\rho_1, \rho_2) < + \infty$. 




\section{Elementary properties of Connes spectral distance}\label{sec3}
In this section, we will review some basic but important properties of Connes spectral distances and the optimal elements. Some of these conclusions have already appeared in other research papers in a similar form. We have listed these simple properties here mainly for the convenience of reference in the following text.
\begin{lemma}\label{d0}\cite{Martinettif}
  For any states $\rho_1, \rho_2$, the Connes spectral distance $d (\rho_1,
  \rho_2) = 0$ if and only if $\rho_1 = \rho_2$. On the other hand, for any
  different states $\rho_1, \rho_2$, there should be $d (\rho_1, \rho_2) > 0$.
\end{lemma}

\begin{proof}
Obviously, for the same states $\rho_1 = \rho_2$,
there is $d (\rho_1, \rho_2) = 0$.

For different states $\rho_1, \rho_2$, there is $\Delta \rho \neq 0$, and $\|
\Delta \rho \| > 0$. If $\|[\mathcal{D}, \pi (\Delta \rho)]\|_{op} = 0$, then
for any number $c > 0$, we always have $\|[\mathcal{D}, \pi (c \Delta
\rho)]\|_{op} = c \|[\mathcal{D}, \pi (\Delta \rho)]\|_{op} = 0$. Therefore,
there is $e_c \equiv c \Delta \rho \in B$ for $\forall c > 0$. But in this
case, we have
\begin{equation}
  d (\rho_1, \rho_2) = \sup_{e \in B} [\mathrm{tr} (\Delta \rho
  \, e)] \geqslant \mathrm{tr} (\Delta \rho \, e_c)
  = c \| \Delta \rho \|^2, \quad \forall c > 0.
\end{equation}
So there should be $d (\rho_1, \rho_2) = + \infty$.

If $\|[\mathcal{D}, \pi (\Delta \rho)]\|_{op} \neq 0$, then one can choose
\begin{equation}
  e = \frac{\Delta \rho}{\|[\mathcal{D}, \pi (\Delta \rho)]\|_{op}},
\end{equation}
and there is $\|[\mathcal{D}, \pi (e)]\|_{op} = 1$. So we have $e \in B$, and
\begin{equation}
  | \mathrm{tr} (\Delta \rho \, e) | = \frac{| \mathrm{tr}
  (\Delta \rho \Delta \rho) |}{\|[\mathcal{D}, \pi (\Delta \rho)]\|_{op}} = \frac{\| \Delta \rho \|^2}{\|[\mathcal{D}, \pi (\Delta
  \rho)]\|_{op}} > 0.
\end{equation}
Therefore, for different states $\rho_1, \rho_2$, there should be some $e \in
B$ which satisfies $| \mathrm{tr} (\Delta \rho \, e) | > 0$. So
we have
\begin{equation}
  d (\rho_1, \rho_2) = \sup_{e \in B} | \mathrm{tr} (\Delta \rho
  \, e) | > 0.
\end{equation}
\end{proof}

\begin{lemma}
  The spectral triples $(\mathcal{A}, \mathcal{H}, \mathcal{D})$ and
  $(\mathcal{A}, \mathcal{H}, \lambda \mathbb{I}+\mathcal{D})$ have the same
  metric, where $\lambda$ is any real number.
\end{lemma}

\begin{proof}
Obviously, there is
\begin{equation}
  B = \{e \in \mathcal{A}: \|[\mathcal{D}, \pi (e)]\|_{op} \leqslant 1\} = \{e
  \in \mathcal{A}: \|[\lambda \mathbb{I}+\mathcal{D}, \pi (e)]\|_{op}
  \leqslant 1\} .
\end{equation}
So the spectral triples $(\mathcal{A}, \mathcal{H}, \mathcal{D})$ and
$(\mathcal{A}, \mathcal{H}, \lambda \mathbb{I}+\mathcal{D})$ have the same
ball conditions and metric.
\end{proof}

So in most of the cases, one can just simply consider the Dirac
operator $\mathcal{D}$ being Hermitian and traceless,
\begin{equation}
  \mathrm{tr} (\mathcal{D}) = 0.
\end{equation}


Furthermore, if we assume that the representation $\pi$ is
\textit{unital}, namely, $\pi (\mathbb{I}) =\mathbb{I}$, we also
have the following result.

\begin{lemma}
  \label{tloe}For any states $\rho_1, \rho_2$, if $e_o$ is the corresponding
  optimal element, then $\lambda \mathbb{I}+ e_o$ is also an optimal element,
  where $\lambda$ is any real number.
\end{lemma}

\begin{proof}
For the optimal element $e_o$, there is
\begin{equation}
  \pi (\lambda \mathbb{I}+ e_o) = \lambda \pi (\mathbb{I}) + \pi (e_o) =
  \lambda \mathbb{I}+ \pi (e_o),
\end{equation}
and
\begin{equation}
  \|[\mathcal{D}, \pi (\lambda \mathbb{I}+ e_o)]\|_{op} = \|[\mathcal{D},
  \lambda \mathbb{I}+ \pi (e_o)]\|_{op} = \|[\mathcal{D}, \pi (e_o)]\|_{op}
  \leqslant 1.
\end{equation}
So $\lambda \mathbb{I}+ e_o \in B$. Since $\mathrm{tr} (\Delta\rho)=\mathrm{tr} (\rho_1-\rho_2)=0$, we also have
\begin{equation}
  \mathrm{tr} [\Delta \rho (\lambda \mathbb{I}+ e_o)] = \mathrm{tr} (\Delta
  \rho \cdot \lambda \mathbb{I}) + \mathrm{tr} (\Delta \rho e_o) = \mathrm{tr}
  (\Delta \rho e_o) = \sup_{e \in B} [\mathrm{tr} (\Delta \rho e)] .
\end{equation}
Therefore, $\lambda \mathbb{I}+ e_o$ is also an optimal element for the states
$\rho_1, \rho_2$.
\end{proof}

So if the representation $\pi$ is unital, one can simply consider the optimal elements $e_o \in B$
being Hermitian and traceless,
\begin{equation}
  \mathrm{tr} (e_o) = 0.
\end{equation}
\begin{lemma}
  Denote the spectral distance between the states $\rho_1, \rho_2$ in the spectral
  triple $(\mathcal{A}, \mathcal{H}, \mathcal{D})$ as $d_{\mathcal{D}}
  (\rho_1, \rho_2)$, and there is
  \begin{equation}
    d_{\lambda \mathcal{D}} (\rho_1, \rho_2) = \frac{1}{| \lambda |}
    d_{\mathcal{D}} (\rho_1, \rho_2),
  \end{equation}
  where $\lambda$ is any real number.
\end{lemma}

\begin{proof}
For any real number $\lambda$, there are
\begin{equation}
  \|[\lambda \mathcal{D}, \pi (e)]\|_{op} = \|[\mathcal{D}, \lambda \pi
  (e)]\|_{op} = \|[\mathcal{D}, \pi (\lambda e)]\|_{op} .
\end{equation}
So if the element $e \in B$ in the spectral triple $(\mathcal{A}, \mathcal{H},
\mathcal{D})$, namely, $\|[\mathcal{D}, \pi (e)]\|_{op} \leqslant 1$, then we
have $\|[\lambda \mathcal{D}, \pi \left( \frac{1}{\lambda} e \right)]\|_{op}
\leqslant 1$, namely, $e' = \frac{1}{\lambda} e \in B'$ in the spectral triple
$(\mathcal{A}, \mathcal{H}, \lambda \mathcal{D})$; vice versa.

So for any states $\rho_1, \rho_2$, there is
\begin{eqnarray}
  d_{\lambda \mathcal{D}} (\rho_1, \rho_2) & = & \sup_{e' \in B'} |
  \mathrm{tr} (\Delta \rho \, e') | \nonumber\\
  & = & \sup_{e \in B} | \mathrm{tr} (\Delta \rho \cdot \frac{1}{\lambda} e)
  | \nonumber\\
  & = & \frac{1}{| \lambda |} \sup_{e \in B} | \mathrm{tr} (\Delta \rho
  \, e) | \nonumber\\
  & = & \frac{1}{| \lambda |} d_{\mathcal{D}} (\rho_1, \rho_2) . 
\end{eqnarray}
\end{proof}

\begin{lemma}\cite{Martinettip}
  Consider any two different states $\rho_1, \rho_2$, and $e_o$ is the
  corresponding optimal element, if the spectral distance is finite, then
  there should be $\|[\mathcal{D}, \pi (e_o)]\|_{op} > 0$.
\end{lemma}

\begin{proof}
Since $\rho_1, \rho_2$ are different states, from Lemma \ref{d0}, there is
$d (\rho_1, \rho_2) > 0$. For the optimal element $e_o$, if there is
$\|[\mathcal{D}, \pi (e_o)]\|_{op} = 0$, and
\begin{equation}
  0 < d (\rho_1, \rho_2) = \sup_{e \in B} [\mathrm{tr} (\Delta \rho
  \, e)] = \mathrm{tr} (\Delta \rho e_o) < + \infty,
\end{equation}
then there should be some real number $t > 1$,
\begin{equation}
  \|[\mathcal{D}, \pi (te_o)]\|_{op} = t \|[\mathcal{D}, \pi (e_o)]\|_{op} = 0
  \leqslant 1.
\end{equation}
So there is $te_o \in B$. But
\begin{equation}
  \mathrm{tr} (\Delta \rho te_o) = t \cdot \mathrm{tr} (\Delta \rho e_o) >
  \mathrm{tr} (\Delta \rho e_o) = \sup_{e \in B} [\mathrm{tr} (\Delta \rho
  \, e)],
\end{equation}
which is impossible. Therefore, if the spectral distance is finite, there should be
$\|[\mathcal{D}, \pi (e_o)]\|_{op} \neq 0$. So we have
$\|[\mathcal{D}, \pi (e_o)]\|_{op} > 0$ for two different states $\rho_1, \rho_2$.
\end{proof}

Obviously, if $\pi (\mathcal{A}) \subseteq \mathcal{C} (\mathcal{D})$, where
  $\mathcal{C} (\mathcal{D})$ is the centralizer of $\mathcal{D}$ in $L
  (\mathcal{H})$, then the Connes spectral distances between any states $\rho_1, \rho_2$ are $d (\rho_1, \rho_2) =
  0$ or $+ \infty$.
The trivial case is that $\mathcal{D}= c\mathbb{I}$. In the present work, we
only study the cases that the Connes spectral distances are finite.

\begin{theorem}\cite{Iochum}
  If the spectral distance of two different states $\rho_1, \rho_2$ is finite,
  then for the corresponding optimal elements $e_o$, there is
  \begin{equation}
    \|[\mathcal{D}, \pi (e_o)]\|_{op} = 1.
  \end{equation}
\end{theorem}

\begin{proof}
Since the spectral distance is finite, for the
optimal element $e_o$, there should be $\|[\mathcal{D}, \pi (e_o)]\|_{op} >
0$. If there are
\begin{equation}
  0 < \|[\mathcal{D}, \pi (e_o)]\|_{op} = t < 1,
\end{equation}
then we also have
\begin{equation}
  \left\| \left[ \mathcal{D}, \pi \left( \frac{e_o}{t} \right) \right]
  \right\|_{op} = \|[\mathcal{D}, \pi (e_o')]\|_{op} = \frac{1}{t}
  \|[\mathcal{D}, \pi (e_o)]\|_{op} = 1.
\end{equation}
So there is $e_o' = \frac{e_o}{t} \in B$, but
\begin{eqnarray}
  \mathrm{tr} (\Delta \rho e_o') & = & \mathrm{tr} (\Delta \rho \frac{e_o}{t})
  \nonumber\\
  & = & \frac{1}{t} \mathrm{tr} (\Delta \rho e_o) \nonumber\\
  & > & \mathrm{tr} (\Delta \rho e_o) \nonumber\\
  & = & \sup_{e \in B} [\mathrm{tr} (\Delta \rho \, e)]
  \nonumber\\
  & = & d (\rho_1, \rho_2) . 
\end{eqnarray}
This is a contradiction. So for any optimal element $e_o$ of different states,
there should be $\|[\mathcal{D}, \pi (e_o)]\|_{op} =1$.
\end{proof}

Since for a given state $\rho$, any element $e_o$ can be consider as the
corresponding optimal element,
\begin{equation}
  d (\rho, \rho) = \mathrm{tr} (\Delta \rho e_o) = 0.
\end{equation}
Therefore, one can just simply choose the optimal element $e_o$ which
satisfies the relation $\|[\mathcal{D}, \pi (e_o)]\|_{op} = 1$.

\begin{definition}
  Let $(\mathcal{A}, \mathcal{H}, \mathcal{D})$ be a spectral triple with
  finite spectral distances, and denote the set of all optimal elements as
  $B_o$. Obviously, there is
  \begin{equation}
    \|[\mathcal{D}, \pi (e)]\|_{op} = 1, \quad \forall e \in B_o .
  \end{equation}
\end{definition}

We also define the set
\begin{equation}
  B_1 = \{e \in \mathcal{A}: \|[\mathcal{D}, \pi (e)]\|_{op} = 1\} .
\end{equation}
Obviously, there is
\begin{equation}
  B_o \subseteq B_1 .
\end{equation}
In the present work, we only consider the spectral triple with finite spectral
distances. So in most of the cases, one can only consider the subset $B_o$ or
$B_1$ instead of $B$.

Actually, we also have the following results.

\begin{lemma}\cite{Martinettip}
  Let $\mathcal{C} (\mathcal{D}) \subset L (\mathcal{H})$ be the centralizer
  of $\mathcal{D}$, if Connes spectral distances in the spectral triples
  $(\mathcal{A}, \mathcal{H}, \mathcal{D})$ are finite, then the subset $\pi
  (\mathcal{A}) \cap \mathcal{C} (\mathcal{D})$ is trivial, namely, $\pi
  (\mathcal{A}) \cap \mathcal{C} (\mathcal{D}) = \{\lambda \mathbb{I}\}$.
\end{lemma}

\begin{proof}
If the subset $\pi (\mathcal{A}) \cap \mathcal{C}
(\mathcal{D})$ is non-trivial, namely, there is some element $e \in
\mathcal{A}$ with $e \neq \lambda \mathbb{I}$ and $\pi (e) \in \mathcal{C}
(\mathcal{D})$. So we have $\|[\mathcal{D}, \pi (e)]\|_{op} = 0$, and $e \in
B$. By the separating property of states, there do exist some states $\rho_1,
\rho_2$ with $\mathrm{tr} (\rho_1  \, e) \neq \mathrm{tr} (\rho_2
\, e)$. Therefore, we have
\begin{equation}
  | \mathrm{tr} (\rho_1  \, e) - \mathrm{tr} (\rho_2 
  \, e) | = | \mathrm{tr} (\Delta \rho \, e) | =
  \delta > 0.
\end{equation}
Obviously, for any real number $k > 0$, we also have $\|[\mathcal{D}, \pi
(ke)]\|_{op} = 0$. So $e' = ke \in B$, but
\begin{equation}
  d (\rho_1, \rho_2) = \sup_{e\in B} | \mathrm{tr} (\Delta \rho e) |
  \geqslant | \mathrm{tr} (\Delta \rho ke) | = k \delta .
\end{equation}
So there should be $d (\rho_1, \rho_2) = + \infty$.

Therefore, if the spectral distances are finite, then the subset $\pi
(\mathcal{A}) \cap \mathcal{C} (\mathcal{D})$ is
trivial.
\end{proof}

\begin{corollary}
  If $\mathcal{A}=\mathbb{M}_n (\mathbb{C})$, $\mathcal{H}=\mathbb{C}^n$, $n
  \geqslant 2$, $\mathcal{D} \in \mathbb{M}_n (\mathbb{C})$, and the
  representation $\pi (a) = a \in \mathbb{M}_n (\mathbb{C})$, then the Connes
  spectral distances in the spectral triple $(\mathcal{A}, \mathcal{H},
  \mathcal{D})$ can not all be finite.
\end{corollary}

\begin{proof}
For any matrix $\mathcal{D}$, the centralizer
$\mathcal{C} (\mathcal{D}) \subseteq \pi (\mathcal{A}) =\mathbb{M}_n
(\mathbb{C})$ is non-trivial. So the subset $\pi (\mathcal{A}) \cap
\mathcal{C} (\mathcal{D})$ is non-trivial, and then Connes spectral distances
in the spectral triple $(\mathcal{A}, \mathcal{H}, \mathcal{D})$ can not all
be finite.
\end{proof}

In Ref.~\cite{Martinettip}, there is also a similar conclusion to the result above.

For example, consider the spectral triple $(\mathbb{M}_2 (\mathbb{C}),
\mathbb{C}^2, \mathcal{D})$ with the Dirac operator
\begin{equation}
  \mathcal{D} = \left(\begin{array}{cc}
    0 & 1\\
    1 & 0
  \end{array}\right),
\end{equation}
and the representation $\pi (a) = a \in \mathbb{M}_2 (\mathbb{C})$. For the element
\begin{equation}
  e_1 = \left(\begin{array}{cc}
    0 & 1\\
    1 & 0
  \end{array}\right) \in \mathbb{M}_2 (\mathbb{C}),
\end{equation}
there is $\|[\mathcal{D}, \pi (e_1)]\|_{op} = 0$, and $k e_1 \in B$ for
any real number $k > 0$.

Consider two states
\begin{equation}
  \rho_1 = |1 \rangle \langle 1| = \left(\begin{array}{cc}
    1 & 0\\
    0 & 0
  \end{array}\right), \quad \rho_2 = | + \rangle \langle + | = \frac{1}{2}
  \left(\begin{array}{cc}
    1 & 1\\
    1 & 1
  \end{array}\right),
\end{equation}
\begin{equation}
  \quad \Delta \rho = \rho_1 - \rho_2 = \frac{1}{2} \left(\begin{array}{cc}
    1 & - 1\\
    - 1 & - 1
  \end{array}\right),
\end{equation}
there is
\begin{equation}
  d (\rho_1, \rho_2) = \sup_{e \in B} | \mathrm{tr} (\Delta \rho e) |
  \geqslant | \mathrm{tr} (\Delta \rho k e_1) | = k | \mathrm{tr} (\Delta \rho
  e_1) | = k.
\end{equation}
So we have $d (\rho_1, \rho_2) = + \infty$.



\section{Unitary invariance of Connes spectral distances of quantum states}\label{sec4}

In quantum physics, most of the properties of physical systems are invariant
under unitary transformations. Connes spectral distances are also essential
physical properties of quantum states. Therefore, it is interesting to study
some types of Connes spectral distances which are unitary invariant. The
unitary invariance of spectral distances can be regarded as a natural property
of a spectral triple. This is also a special kind of isometry of the state space being studied in some other literature, such as Ref.~\cite{Martinettip}.

\begin{definition}
  Assume that the spectral distances in the spectral triple $(\mathcal{A},
  \mathcal{H}, \mathcal{D})$ are all finite. For any states $\rho_1, \rho_2
  \in \mathcal{A}$ and any unitary operator $U \in \mathcal{U} (\mathcal{A})$,
  denote $\rho_1' = U \rho_1 U^{\dag}$, $\rho_2' = U \rho_2 U^{\dag}$. If the
  spectral distances satisfy
  \begin{equation}
    d (\rho_1', \rho_2') = d (\rho_1, \rho_2),
  \end{equation}
  we say that the spectral distances in the spectral triple $(\mathcal{A},
  \mathcal{H}, \mathcal{D})$ are \textit{unitary invariant}.
\end{definition}

Obviously, if $\mathcal{A}$ is commutative, we have $\rho_i' = \rho_i$, and
the spectral distances are unitary invariant.

First, let us consider the sets $B$ and $B_o, B_1$ under unitary
transformations. Usually, the set $B$ is not necessarily closed under unitary
conjugations, namely, if $e \in B$, then $UeU^{\dag} \in B$, where $U \in
\mathcal{U} (\mathcal{A})$ is any unitary operator. Similarly, the sets $B_o,
B_1$ are also not necessarily closed under unitary conjugations.

\begin{lemma}
  Consider a spectral triple $(\mathcal{A}, \mathcal{H}, \mathcal{D})$ with
  finite Connes spectral distances, if the optimal element set $B_o$ is closed
  under unitary conjugations, then the Connes spectral distances of the states
  are unitary invariant.
\end{lemma}

\begin{proof}
If the optimal element set $B_o$ is closed under
unitary conjunctions, namely, $e_o \in B_o$ if and only if $Ue_o U^{\dag} \in
B_o$, where $U \in \mathcal{U} (\mathcal{A})$ is any unitary operator, then
for any state $\rho_1, \rho_2$, and $\rho_i' = U \rho_i U^{\dag}$, $\Delta
\rho' = U \Delta \rho U^{\dag}$, there are
\begin{eqnarray}
  d (\rho_1, \rho_2) & = & \sup_{e \in B_o} | \mathrm{tr} (\Delta \rho e) |
  \nonumber\\
  & = & \sup_{UeU^{\dag} \in B_o} | \mathrm{tr} (\Delta \rho e) | \nonumber\\
  & = & \sup_{UeU^{\dag} \in B_o} | \mathrm{tr} (U \Delta \rho U^{\dag} \cdot
  UeU^{\dag}) | \nonumber\\
  & = & \sup_{e' \in B_o} | \mathrm{tr} (\Delta \rho' e') | \nonumber\\
  & = & d (\rho_1', \rho_2') . 
\end{eqnarray}
So the Connes spectral distances are unitary
invariant.
\end{proof}

Similarly, one can prove the following result.

\begin{lemma}
  If the set $B_1$ (or $B$) is closed under unitary conjugations, then the
  Connes spectral distances are unitary invariant.
\end{lemma}

\begin{corollary}
  If the set $B_1$ (or $B$) is closed under unitary conjugations, so is the
  optimal element set $B_o$.
\end{corollary}

\begin{proof}Since the set $B_1$ is closed under unitary
conjugation, the Connes spectral distances are unitary invariant. For any
optimal element $e_o \in B_o$ and any unitary operator $U$, there are some
states $\rho_1, \rho_2$ corresponding to $e_o$, and $\rho_i' = U \rho_i
U^{\dag}$,
\begin{eqnarray}
  \mathrm{tr} (\Delta \rho e_o) & = & d (\rho_1, \rho_2) \nonumber\\
  & = & \mathrm{tr} (U \Delta \rho U^{\dag} \cdot Ue_o U^{\dag}) \nonumber\\
  & = & \mathrm{tr} (\Delta \rho' e_o') . \nonumber\\
  & = & d (\rho_1', \rho_2') \nonumber\\
  & = & \sup_{e \in B_1} [\mathrm{tr} (\Delta \rho' e)] . 
\end{eqnarray}
Since $e_o \in B_o \subseteq B_1$, we have $Ue_o U^{\dag} = e_o' \in B_1$. So
$Ue_o U^{\dag} = e_o'$ is the corresponding optimal element of $\rho_1',
\rho_2'$. This means that $Ue_o U^{\dag} = e_o' \in B_o$. Therefore, the
optimal element set $B_o$ is closed under unitary
conjunctions.
\end{proof}

\begin{theorem}\label{lu}
  If the operator norms $\|[\mathcal{D}, \pi (e)]\|_{op}$ are invariant under
  any unitary transformations of any elements $e$, namely, for any unitary
  operator $U \in \mathcal{U} (\mathcal{A})$, there is $\|[\mathcal{D}, \pi
  (e)]\|_{op} = \|[\mathcal{D}, \pi (UeU^{\dag})]\|_{op}$, then the Connes
  spectral distances are unitary invariant.
\end{theorem}

\begin{proof}If $\|[\mathcal{D}, \pi (e)]\|_{op}$ are invariant
under any unitary transformations of the elements $e$, then the sets $B$,
$B_1$ and $B_o$ are all closed under unitary conjugations. Therefore, the
Connes spectral distances are unitary invariant.
\end{proof}

This also means that the Lipschitz seminorm $L (a) = \|[\mathcal{D}, \pi
(a)]\|_{op}$ is invariant under unitary conjugation.
One simple case is $[\mathcal{D}, \pi (U)]=0$, and this is also studied in Ref.~\cite{Martinettip}.

Furthermore, if the representation $\pi$ is \textit{unital}, namely,
$\pi (\mathbb{I}) =\mathbb{I}$, we have $\pi (UU^{\dag}) = \pi (\mathbb{I}) =
\pi (U) \pi (U^{\dag}) = \pi (U) \pi (U)^{\dag} =\mathbb{I}$. In this case,
$\pi (U)$ is also unitary. We have the following result.

\begin{corollary}
  If for any unitary operator $U$, there is $\|[\pi (U)\mathcal{D} \pi
  (U^{\dag}), \pi (e)]\|_{op} = \|[\mathcal{D}, \pi (e)]\|_{op}$, then the
  Connes spectral distances are unitary invariant.
\end{corollary}

\begin{proof}
For any unitary operator $U$, there are
\begin{eqnarray}
  \|[\mathcal{D}, \pi (e)]\|_{op} & = & \|[\pi (U)\mathcal{D} \pi (U^{\dag}),
  \pi (e)]\|_{op} \nonumber\\
  & = & \| \pi (U^{\dag}) [\pi (U)\mathcal{D} \pi (U^{\dag}), \pi (e)] \pi
  (U)\|_{op} \nonumber\\
  & = & \|[\mathcal{D}, \pi (U^{\dag}) \pi (e) \pi (U)]\|_{op} \nonumber\\
  & = & \|[\mathcal{D}, \pi (U^{\dag} eU)]\|_{op} . 
\end{eqnarray}
So $\|[\mathcal{D}, \pi (e)]\|_{op}$ is invariant under any unitary
transformations of $e$, from Theorem \ref{lu}, the Connes spectral distances are unitary
invariant.
\end{proof}

\begin{corollary}
  If for any unitary operators $U, U' \in \mathcal{U} (\mathcal{A})$, there is
  $\|[\mathcal{D}, \pi (e)]\|_{op} = \|[\pi (U)\mathcal{D} \pi (U^{\dag}), \pi
  (U' {eU'}^{\dag})]\|_{op}$, then the Connes spectral distances are unitary
  invariant.
\end{corollary}


\begin{theorem}
  \label{udu}The Connes spectral distances in the spectral triple
  $(\mathcal{A}, \mathcal{H}, \mathcal{D})$ are unitary invariant if and only
  if the spectral triples $(\mathcal{A}, \mathcal{H}, \mathcal{D})$ and
  $(\mathcal{A}, \mathcal{H}, \pi (U)\mathcal{D} \pi (U^{\dag}))$ have the
  same metric for any unitary operator $U \in \mathcal{U} (\mathcal{A})$.
\end{theorem}

\begin{proof}
Denote
$\mathcal{D}' = \pi (U) \mathcal{D} \pi (U^{\dag})$,
and $B' \equiv \{e \in \mathcal{A}: \|[\mathcal{D}', \pi (e)]\|_{op}
\leqslant 1\}$ .
Obviously, for the given unitary operator $U$, if $e \in B$ in the
spectral triple $(\mathcal{A}, \mathcal{H}, \mathcal{D})$, then there is $e' =
UeU^{\dag} \in B'$ in $(\mathcal{A}, \mathcal{H}, \mathcal{D}')$; vice
versa.

If the Connes spectral distances in the spectral
triple $(\mathcal{A}, \mathcal{H}, \mathcal{D})$ are unitary invariant, then
for any states $\rho_1, \rho_2$, and $\rho_i' = U^{\dag} \rho_i U$, there is
\begin{eqnarray}
  d_{\mathcal{D}} (\rho_1, \rho_2) & = & \sup_{e \in B} | \mathrm{tr} (\Delta
  \rho e) | = d_{\mathcal{D}} (\rho_1', \rho_2') \nonumber\\
  & = & \sup_{e \in B} | \mathrm{tr} (\Delta \rho' e) | = \sup_{e \in B} |
  \mathrm{tr} (U^{\dag} \Delta \rho U e) | \nonumber\\
  & = & \sup_{e \in B} | \mathrm{tr} (\Delta \rho U eU^{\dag}) | = \sup_{e
  \in B} | \mathrm{tr} (\Delta \rho e') | \nonumber\\
  & = & \sup_{e' \in B'} | \mathrm{tr} (\Delta \rho e') | \nonumber\\
  & = & d_{\mathcal{D}'} (\rho_1, \rho_2) . 
\end{eqnarray}
Therefore, any two states $\rho_1, \rho_2\in\mathcal{A}$ have the same spectral distances in
the spectral triples $(\mathcal{A}, \mathcal{H}, \mathcal{D})$ and
$(\mathcal{A}, \mathcal{H}, \pi (U)\mathcal{D} \pi (U^{\dag}))$. So
$(\mathcal{A}, \mathcal{H}, \mathcal{D})$ and $(\mathcal{A}, \mathcal{H}, \pi
(U)\mathcal{D} \pi (U^{\dag}))$ have the same metric.

On the other hand, for any two states $\rho_1, \rho_2\in\mathcal{A}$, if they have the same
spectral distances in these two spectral triples, then we have
\begin{eqnarray}
  d_{\mathcal{D}} (\rho_1, \rho_2) & = & \sup_{e \in B} | \mathrm{tr} (\Delta
  \rho \, e) | \nonumber\\
  & = & d_{\mathcal{D}'} (\rho_1, \rho_2) \nonumber\\
  & = & \sup_{e' \in B'} | \mathrm{tr} (\Delta \rho e') | \nonumber\\
  & = & \sup_{e \in B} | \mathrm{tr} (\Delta \rho UeU^{\dag}) | \nonumber\\
  & = & \sup_{e \in B} | \mathrm{tr} (U^{\dag} \Delta \rho U \,
  e) | \nonumber\\
  & = & d_{\mathcal{D}} (\rho_1', \rho_2') . 
\end{eqnarray}
Therefore, the Connes spectral distances between any states $\rho_1, \rho_2$
are unitary invariant.
\end{proof}

\begin{corollary}
  The Connes spectral distances in the spectral triple $(\mathcal{A},
  \mathcal{H}, \mathcal{D})$ are unitary invariant if and only if the Connes
  spectral distances in $(\mathcal{A}, \mathcal{H}, \pi (U)\mathcal{D} \pi
  (U^{\dag}))$ are unitary invariant, where $U \in \mathcal{U} (\mathcal{A})$
  is any unitary operator.
\end{corollary}

There are also some other articles have studied the question of isometries for some types of spectral triples, such as Ref.~\cite{Bassi}.




\section{Spectral triples with quantum trace distances}\label{sec5}

Quantum trace distance is one of the most useful distance measures between
quantum states in quantum information sciences. It is also unitary invariant.
So it is also very significant to study spectral triples with Connes spectral
distances which equal quantum trace distances.

\begin{definition}
  \cite{Nielsen} For two states $\rho_1, \rho_2$, the quantum trace distance
  is defined as
  \begin{equation}
    d_{\mathrm{tr}} (\rho_1, \rho_2) = \| \rho_1 - \rho_2 \|_1 = \| \Delta
    \rho \|_1 = \mathrm{tr} \sqrt{\Delta \rho^{\dag} \Delta \rho},
  \end{equation}
  where the trace norm $\|A\|_1 := \mathrm{tr} \sqrt{A^{\dag} A}$. Note that,
  here we have ignored the coefficient $\frac{1}{2}$ for simplicity.
\end{definition}

\begin{theorem}\label{dtr}
  \cite{Nielsen} For two states $\rho_1, \rho_2$, the quantum trace distance
  is equal to
  \begin{equation}
    d_{\mathrm{tr}} (\rho_1, \rho_2) = \sup_{-\mathbb{I} \leqslant P \leqslant
    \mathbb{I}} \left| \mathrm{tr} [P (\rho_1 - \rho_2)] \right| = \sup_{\|P\|_{op}
    \leqslant 1} | \mathrm{tr} (P \Delta \rho) |,
  \end{equation}
  where $P$ are Hermitian operators.
\end{theorem}

\begin{theorem}
  Consider the spectral triple $(\mathcal{A}, \mathcal{H}, \mathcal{D})$ with
  $\mathcal{A}=\mathbb{M}_n (\mathbb{C})$, there are some
  representation $\pi$ of $\mathcal{A}$ on some Hilbert space $\mathcal{H}$,
  and some Dirac operator $\mathcal{D} \in L (\mathcal{H})$ which satisfy the
  relation $\|[\mathcal{D}, \pi (e)]\|_{op} = \|e\|_{op}$ for any element $e
  \in \mathcal{A}$.
\end{theorem}

To proof this theorem, we will construct a concrete example.

\begin{proof}
For $\mathcal{A}=\mathbb{M}_n (\mathbb{C})$, one can
choose the Hilbert space $\mathcal{H}=\mathbb{C}^n \oplus \mathbb{C}^n$, and
the representation $\pi (a) = \left( \begin{array}{cc}
  a & 0\\
  0 & 0
\end{array} \right)$ for $a \in \mathcal{A}$. The Dirac operator is
\begin{equation}
  D = \left( \begin{array}{cc}
    0 & \mathbb{I}_n\\
    \mathbb{I}_n & 0
  \end{array} \right) .
\end{equation}
The commutator $[D, \pi (a)]$ is:
\begin{equation}
  [D, \pi (a)] = \left( \begin{array}{cc}
    0 & - a\\
    a & 0
  \end{array} \right),
\end{equation}
and
\begin{equation}
  [D, \pi (a)]^{\dag} [D, \pi (a)] = \left( \begin{array}{cc}
    a^{\dag} a & 0\\
    0 & a^{\dag} a
  \end{array} \right) .
\end{equation}
So the eigenvalues of $[D, \pi (a)]^{\dag} [D, \pi (a)]$ are exactly the
eigenvalues of $a^{\dag} a$. Therefore, we have
\begin{equation}
  \|[D, \pi (a)]\|_{op} = \|a\|_{op} .
\end{equation}
\end{proof}

Note that the representation above is not unital. In this case,
the Lipschitz seminorm is just the normal operator norm. It is unitary invariant, and then the Connes spectral distances are unitary invariant.

\begin{theorem}
  \label{t6}In the spectral triple $(\mathcal{A}, \mathcal{H}, \mathcal{D})$,
  if the Dirac operator $\mathcal{D}$ satisfies the relation $\|[\mathcal{D},
  \pi (e)]\|_{op} = \|e\|_{op}$ for any element $e \in \mathcal{A}$, then the
  ball condition becomes
  \begin{equation}
    B = \{e \in \mathcal{A}: \|e\|_{op} \leqslant 1\},
  \end{equation}
  and the Connes spectral distances are equal to the quantum trace distances,
  \begin{equation}
    d (\rho_1, \rho_2) = d_{\mathrm{tr}} (\rho_1, \rho_2) = \| \rho_1 - \rho_2
    \|_1 .
  \end{equation}
\end{theorem}

\begin{proof}
From Theorem \ref{dtr}, in this case, the Connes spectral distance is
\begin{eqnarray}
  d (\rho_1, \rho_2) & = & \sup_{\|e\|_{op} \leqslant 1} | \mathrm{tr} (\Delta
  \rho \, e) | \nonumber\\
  & = & d_{\mathrm{tr}} (\rho_1, \rho_2) \nonumber\\
  & = & \| \rho_1 - \rho_2 \|_1 .
\end{eqnarray}
\end{proof}

Since the Hermitian operator $\Delta \rho = \rho_1 - \rho_2$ can be
represented by
\begin{equation}
  \Delta \rho = \sum_i a_i |i \rangle \langle i|, \quad \sum_i a_i = 0, \quad
  a_i \neq 0,
\end{equation}
we have
\begin{equation}
  d (\rho_1, \rho_2) = d_{\mathrm{tr}} (\rho_1, \rho_2) =\| \rho_1 - \rho_2 \|_1 = \mathrm{tr}
  \sqrt{\Delta \rho^{\dag} \Delta \rho} = \sum_i |a_i |,
\end{equation}
and the optimal element $e_o$ can be chosen as
\begin{equation}
  e_o = \sum_i \frac{a_i}{|a_i |} |i \rangle \langle i|. \label{oe}
\end{equation}
One can easily verify that, there are
\begin{equation}
  \|e_o \|_{op} = 1, \quad \mathrm{tr} (\Delta \rho \, e_o) =
  \sum_i |a_i |=d (\rho_1, \rho_2) .
\end{equation}
It is easy to see that, $e_o$ (\ref{oe}) is the only possible Hermitian
optimal element.


Furthermore, if the representation $\pi$ is \textit{unital}, from Lemma (\ref{tloe}),
one can choose all optimal elements $e_o$ to be traceless. From the result
(\ref{oe}), one can find that $\mathcal{A}=\mathbb{M}_2 (\mathbb{C})$ is the
only possible case in which all optimal elements can be traceless.

\begin{theorem}
  Consider the spectral triple $(\mathcal{A}, \mathcal{H}, \mathcal{D})$ with
  $\mathcal{A}=\mathbb{M}_n (\mathbb{C})$, and some unital
  representation $\pi$ of $\mathcal{A}$ on the Hilbert space $\mathcal{H}$,
  if the Dirac operator $\mathcal{D} \in L (\mathcal{H})$ satisfies the
  relation $\|[\mathcal{D}, \pi (e_o)]\|_{op} = \|e_o \|_{op}$ for any optimal
  elements $e_o$ of Connes spectral distances, then there should be
  $\mathcal{A}=\mathbb{M}_2 (\mathbb{C})$.
\end{theorem}
This just corresponds to the case of one-qubit states.

There are also many examples with the algebra $\mathcal{A}$ not being some matrix algebra $\mathbb{M}_n (\mathbb{C})$.
For example, one of the simplest case is the two point space.
Define the spectral triple $(\mathcal{A}, \mathcal{H}, \mathcal{D})$ with $\mathcal{A}\cong\mathbb{C}^2$ (diagonal $2 \times 2$ matrices),
$\mathcal{H}=\mathbb{C}^2$ and $D = \frac{1}{2} \left(\begin{array}{cc}
	0 & 1\\
	1 & 0
\end{array}\right)$, and $\pi(e)=e$ for $e\in\mathcal{A}$. The representation $\pi$ is also linear and unital. For any traceless Hermitian
element $e = \left(\begin{array}{cc}
	a & 0\\
	0 & - a
\end{array}\right) \in \mathcal{A}$, the Lipschitz seminorm is
\begin{equation}
\| [D, e] \|_{o p} = \left| \text{$a$} \right| = \|e\|_{op} .
\end{equation}
For any two states
\begin{equation}
\rho_1 = \left(\begin{array}{cc}
	p & 0\\
	0 & 1 - p
\end{array}\right),\qquad
\rho_2 = \left(\begin{array}{cc}
	q & 0\\
	0 & 1 - q
\end{array}\right),
\end{equation}
where $0\leqslant p,q\leqslant1$, and
\begin{equation}
\Delta\rho =\rho_1-\rho_2= \left(\begin{array}{cc}
p-q & 0\\
0 & q - p
\end{array}\right).
\end{equation}
The Connes distance is:
\begin{equation}
d (\rho_1,\rho_2) = \sup_{e \in B}[ \mathrm{tr} (\Delta \rho
\, e)] = \sup_{|a| \le 1} 2 (p - q) a = 2 |p - q| =
d_{\mathrm{tr}} (\rho_1,\rho_2) .
\end{equation}
Therefore, the spectral distance exactly matches the trace distance. The algebra $\mathcal{A}$ is commutative, so the spectral distances are unitary invariant.

In Ref.~\cite{Kumar1}, the authors also constructed a concrete spectral triple related to the trace distance.




\section{Connes spectral distances of one-qubit states}\label{sec6}
In this section, we will study some concrete examples for one-qubit states. These spectral triples can help us study the geometric structure of quantum states.

Similar to the construction in Ref. \cite{LX}, one can define a spectral
triple as follows.
\begin{equation}
	\mathcal{A}= \{ |i \rangle
	\langle j| \}\cong\mathbb{M}_2 (\mathbb{C}) ,\qquad \mathcal{H}_4  =\mathbb{C}^2 \otimes F\cong\mathbb{C}^4,\qquad F = \{
|0 \rangle, |1 \rangle \},
\end{equation}
and $\pi (a) =\mathbb{I}_2 \otimes a = \left(
\begin{array}{cc}
  a & 0\\
  0 & a
\end{array} \right)$, $\forall a \in \mathcal{A}$. The Dirac operator is
\begin{equation}\label{dd4}
  \mathcal{D}_4 = \frac{1}{4}  \sum_{i = 1}^3 \sigma_i \otimes \sigma_i =
  \frac{1}{4}  (\sigma_1 \otimes \sigma_1 + \sigma_2 \otimes \sigma_2 +
  \sigma_3 \otimes \sigma_3),
\end{equation}
where $\sigma_i$ are the Pauli matrices. In this case, the representation
$\pi$ is linear and unital.

In general, the density matrix $\rho$ for a qubit can be expressed as
\cite{Nielsen}
\begin{equation}
  \label{rr} \rho = \frac{I + \vec{r} \cdot \vec{\sigma}}{2},
\end{equation}
where the real vector $\vec{r}$ are the so-called Bloch vector, $|
\vec{r} | \leqslant 1$, and $\vec{\sigma} = (\sigma_1, \sigma_2, \sigma_3)$.

\begin{lemma}
  In the spectral triple $(\mathbb{M}_2 (\mathbb{C}), \mathbb{C}^4,
  \mathcal{D}_4)$, for any traceless Hermitian matrix $e \in \mathbb{M}_2
  (\mathbb{C})$, there is $\|[\mathcal{D}_4, \pi (e)]\|_{op} = \|e\|_{op}$.
\end{lemma}

\begin{proof}
	Any $2 \times 2$ traceless Hermitian matrix $e$ can be
expressed by the Pauli matrices as
\begin{equation}
	e = \vec{e} \cdot \vec{\sigma} = e_1 \sigma_1 + e_2 \sigma_2 + e_3 \sigma_3,
\end{equation}
where $\vec{e} = (e_1, e_2, e_3)$ is a real vector.
Using the following formulas of the Pauli matrices
\begin{equation}
	\sigma_i \sigma_j = \delta_{i j} \mathbb{I}+ \text{i} \varepsilon_{i j k}
	\sigma_k,
\end{equation}
\begin{equation}
	(\vec{a} \cdot \vec{\sigma})  (\vec{b} \cdot \vec{\sigma}) = (\vec{a} \cdot
	\vec{b}) \mathbb{I}_2 + \text{i}  (\vec{a} \times \vec{b}) \cdot
	\vec{\sigma},
\end{equation}
\begin{equation}
	[\vec{a} \cdot \vec{\sigma}, \vec{b} \cdot \vec{\sigma}] = 2 \text{i}
	(\vec{a} \times \vec{b}) \cdot \vec{\sigma},
\end{equation}
after some straightforward calculations, one can obtain
\begin{equation}
	e^{\dag} e = | \vec{e}  |^2 \cdot \mathbb{I}_2,
\end{equation}
\begin{equation}
	[\mathcal{D}_4, \mathbb{I}_2 \otimes e]^{\dag} [\mathcal{D}_4, \mathbb{I}_2
	\otimes e] = \frac{1}{2} (| \vec{e}  |^2 \cdot \mathbb{I}_4 - e \otimes e).
\end{equation}
The eigenvalues of $e = \vec{e} \cdot \vec{\sigma}$ are $\{ | \vec{e} |,- | \vec{e} | \}$.
Since the eigenvalues of a tensor product of two matrices $A \otimes B$ are the
products of the eigenvalues of $A$ and the eigenvalues of $B$, the eigenvalues of $e \otimes e$ are $\{| \vec{e} |^2, | \vec{e} |^2, - | \vec{e} |^2,
- | \vec{e} |^2 \}$.
So the eigenvalues of $\frac{1}{2}  (| \vec{e} |^2 \mathbb{I}_4 - e \otimes e)$
are $\{0, 0, | \vec{e} |^2, | \vec{e} |^2 \}$.


Therefore, we have
\begin{equation}
  \|[\mathcal{D}_4, \pi (e)]\|_{op} = | \vec{e} | = \|e\|_{op} .
\end{equation}
\end{proof}


Consider any one-qubit states $\rho_1, \rho_2$, and the corresponding Bloch
vectors are $\vec{r}_1$, $\vec{r}_2$, respectively.
From the above lemma and Theorem (\ref{t6}), one can obtain the Connes spectral distance,
\begin{equation}
	d (\rho_1, \rho_2) = d_{\mathrm{tr}} (\rho_1, \rho_2) = | \Delta \vec{r} | = \mathrm{tr} (\Delta
	\rho e_o) .
\end{equation}
In this case, there is
\begin{equation}
	\Delta \rho =\rho_1-\rho_2= \frac{1}{2} (\vec{r}_1-\vec{r}_2) \cdot \vec{\sigma}= \frac{1}{2} \Delta \vec{r} \cdot \vec{\sigma},
\end{equation}
and the traceless optimal element can be chosen as
\begin{equation}
	e_o = \frac{\Delta \vec{r}}{| \Delta \vec{r} |} \cdot \vec{\sigma} =
	\frac{\Delta \rho}{\| \Delta \rho \|_{op}}.
\end{equation}
The optimal element set $B_o$ is
\begin{equation}
	B_o = \{ \vec{n} \cdot \vec{\sigma} : | \vec{n} | = 1\} .
\end{equation}
Since a unitary transformation of $\vec{n} \cdot \vec{\sigma}$ is just some
rotation of the unit vector $\vec{n}$, the set $B_o$ is unitary invariant.

Therefore, we have the following result.
\begin{corollary}
  In the spectral triple $(\mathbb{M}_2 (\mathbb{C}), \mathbb{C}^4,
  \mathcal{D}_4)$, the Connes spectral distance between any one-qubit states
  $\rho_1, \rho_2$ is equal to their quantum trace distance and also the Euclidean
  distance between the Bloch vectors,
  \begin{equation}
    d (\rho_1, \rho_2) = d_{\mathrm{tr}} (\rho_1, \rho_2) = | \vec{r}_1 - \vec{r}_2 |,
  \end{equation}
  where $\vec{r}_1$, $\vec{r}_2$ are the corresponding Bloch vectors of
  $\rho_1, \rho_2$, respectively. The corresponding optimal element $e_o$ is
  \begin{equation}
    e_o = \frac{\Delta \vec{r}}{| \Delta \vec{r} |} \cdot \vec{\sigma}= \frac{\Delta \rho}{\| \Delta \rho \|_{op}} .
  \end{equation}
\end{corollary}

It is easy to see that, for any unitary matrix $U \in \mathbb{M}_2
(\mathbb{C})$, the Dirac operator $\mathcal{D}_4' = \pi (U) \mathcal{D}_4 \pi
(U^{\dag}) = (\mathbb{I}_2 \otimes U) \mathcal{D}_4  (\mathbb{I}_2 \otimes
U^{\dag})$ also satisfies the property $\|[\mathcal{D}'_4, \pi (e)]\|_{op} =
\|e\|_{op}$ for any traceless Hermitian matrix $e$,
\begin{eqnarray}
  \|[\mathcal{D}'_4, \pi (e)]\|_{op} & = & \|[\pi (U)\mathcal{D}_4 \pi
  (U^{\dag}), \pi (e)]\|_{op}\nonumber\\
  & = & \|[\mathcal{D}_4, \pi (U^{\dag}) \pi (e) \pi (U)]\|_{op}\nonumber\\
  & = & \|[\mathcal{D}_4, \pi (U^{\dag} e U)]\|_{op}\nonumber\\
  & = & \|U^{\dag} e U\|_{op}\nonumber\\
  & = & \|e\|_{op} .
\end{eqnarray}
For example, there are $2 \times 2$ unitary matrices $U_+, U_-$ which can
permutate the Pauli matrices,
\begin{equation}
  U_+ = \frac{1}{2} \left( \begin{array}{cc}
    1 - i & - 1 - i\\
    1 - i & 1 + i
  \end{array} \right), \quad U_- = \frac{1}{2} \left( \begin{array}{cc}
    1 + i & 1 + i\\
    - 1 + i & 1 - i
  \end{array} \right),
\end{equation}
and
\begin{equation}
  U_+ \sigma_1 U_+^{\dag} = \sigma_2, \quad U_+ \sigma_2 U_+^{\dag} =
  \sigma_3, \quad U_+ \sigma_3 U_+^{\dag} = \sigma_1 ;
\end{equation}
\begin{equation}
  U_- \sigma_3 U_-^{\dag} = \sigma_2, \quad U_- \sigma_2 U_-^{\dag} =
  \sigma_1, \quad U_- \sigma_1 U_-^{\dag} = \sigma_3 .
\end{equation}
Therefore, using the Dirac operator $\mathcal{D}_4$ (\ref{dd4}) and the above unitary matrices $U_{\pm}$, one can obtain the following result.
\begin{lemma}
  For any traceless Hermitian matrix $e \in \mathbb{M}_2 (\mathbb{C})$, the
  Dirac operators
  \begin{equation}\label{d4p}
    \mathcal{D}_4' = \frac{1}{4}  \sum_{i = 1}^3 \pm \sigma_i \otimes
    \sigma_{i'}
  \end{equation}
  satisfy the relation $\|[\mathcal{D}_4', \pi (e)]\|_{op} = \|e\|_{op}$,
  where $\sigma_i$ are the Pauli matrices, and $i, i'$ are any permutations of
  $\{1, 2, 3\}$.
\end{lemma}

\begin{theorem}
  Consider the Dirac operator,
  \begin{equation}
    \mathcal{D}_{4^n} = \frac{1}{4}  \sum_{i = 1}^3 \sigma_{i_n, i_n'}
    \otimes \cdots \otimes \sigma_{i_1, i_1'},
  \end{equation}
  where $i_k, i_k'$ are any permutations of $\{1, 2, 3\}$, and $\sigma_{i_k,
  i_k'} = \pm \sigma_{i_k} \otimes \sigma_{i_k'}$. Then for any $2 \times 2$
  traceless Hermitian matrix $e \in \mathbb{M}_2 (\mathbb{C})$, we have
  \begin{equation}
    \|[\mathcal{D}_{4^n}, \pi_n (e)]\|_{op} = \|e\|_{op},
  \end{equation}
  where $\pi_n (e) =\mathbb{I}_{4^{n - 1}} \otimes \mathbb{I}_2 \otimes e$.
\end{theorem}

\begin{proof}
	When $n = 1$, the Dirac operator $\mathcal{D}_{4}$ is just $\mathcal{D}_4'$ (\ref{d4p}).

When $n = 2$, we have
\begin{equation}
	\mathcal{D}_{16} = \frac{1}{4}  \sum_{i = 1}^3 \sigma_{i_2, i_2'} \otimes
	\sigma_{i_1, i_1'},
\end{equation}
and
\begin{eqnarray}
	\|[\mathcal{D}_{16}, \pi_2 (e)]\|_{op} & = & \left\| \left[ \frac{1}{4} 
	\sum_{i = 1}^3 \sigma_{i_2, i_2'} \otimes \sigma_{i_1, i_1'}, \mathbb{I}_4
	\otimes \pi_1 (e) \right] \right\|_{op} \nonumber\\
	& = & \frac{1}{4} \left\| \sum_{i = 1}^3 \sigma_{i_2, i_2'} \otimes
	[\sigma_{i_1, i_1'}, \pi_1 (e)] \right\|_{op} . 
\end{eqnarray}
Using the unitary matrices $U_{\pm}$, one can permutate the Pauli matrices. So
one can only consider the following simplest case,
\begin{equation}
	\|[\mathcal{D}_{16}, \pi_2 (e)]\|_{op} = \frac{1}{4} \left\| \sum_{i = 1}^3
	\sigma_{i, i} \otimes [\sigma_{i, i}, \pi_1 (e)] \right\|_{op} .
\end{equation}
There is the following $4 \times 4$ unitary matrix
\begin{equation}
	U = \frac{1}{\sqrt{2}} \left( \begin{array}{cccc}
		1 & 0 & 0 & - 1\\
		0 & 1 & - 1 & 0\\
		0 & 1 & 1 & 0\\
		1 & 0 & 0 & 1
	\end{array} \right),
\end{equation}
and we have
\begin{eqnarray}
	\|[\mathcal{D}_{16}, \pi_2 (e)]\|_{op} 
	& = & \left\| (U \otimes \mathbb{I}_4) \left( \frac{1}{4}  \sum_{i = 1}^3
	\sigma_{i, i} \otimes [\sigma_{i, i}, \pi_1 (e)] \right) (U^{\dag} \otimes
	\mathbb{I}_4) \right\|_{op} \\
	& = & \left\| \left( \begin{array}{cccc}
		{}[D_1, \pi_1 (e)] & 0 & 0 & 0\\
		0 & [D_2, \pi_1 (e)] & 0 & 0\\
		0 & 0 & [D_3, \pi_1 (e)] & 0\\
		0 & 0 & 0 & [D_4, \pi_1 (e)]
	\end{array} \right) \right\|_{op}, \nonumber
\end{eqnarray}
where
\begin{equation}
	D_k = \frac{1}{4}  \sum_{i = 1}^3 \sigma_{i, i} = \frac{1}{4}  \sum_{i =
		1}^3 \pm \sigma_i \otimes \sigma_i, \quad k = 1, 2, 3, 4,
\end{equation}
and different $D_k$ has different signs of some terms, but they all have the
same relation
\begin{equation}
	\|[D_k, \pi_1 (e)]\|_{op} =\|[\mathcal{D}_4, \pi_1 (e)]\|_{op}= \|e\|_{op}, \quad k = 1, 2, 3, 4 .
\end{equation}
So we have
\begin{equation}
	\|[\mathcal{D}_{16}, \pi_2 (e)]\|_{op} = \|[\mathcal{D}_4, \pi_1 (e)]\|_{op}
	= \|e\|_{op} . 
\end{equation}
Using the similar methods, one can easily find that, for any $n \geqslant 1$, there is
\begin{equation}
	\|[\mathcal{D}_{4^n}, \pi_n (e)]\|_{op} = \|[\mathcal{D}_4, \pi_1 (e)]\|_{op}
	= \|e\|_{op} .
\end{equation}
\end{proof}

For example, one can check that, with the following Dirac operator
\begin{equation}
	\mathcal{D}_{16} = \frac{1}{4}  \sum_{i = 1}^3 \sigma_i \otimes \sigma_i
	\otimes \sigma_i \otimes \sigma_i, \quad {\pi_2}  (e) =\mathbb{I}_8 \otimes
	e,
\end{equation}
for any $2 \times 2$
traceless Hermitian matrix $e \in \mathbb{M}_2 (\mathbb{C})$, there is
\begin{equation}
	\|[\mathcal{D}_{16}, \pi_2 (e)]\|_{op} = \|e\|_{op} .
\end{equation}


\begin{lemma}
  Consider a matrix $A = \sum_{i, j} S_i \otimes T_j$, where $S_i$ are some $k
  \times k$ matrices, and $T_j$ are some $l \times l$ matrices, and define a
  matrix
  \begin{equation}
    A' = \sum_{i, j} S_i \otimes \tilde{M} \otimes T_j,
  \end{equation}
  where $M$ is any non-zero $n \times n$ matrix, and $\tilde{M} =
  \frac{M}{\|M\|_{op}}$, then there is $\|A\|_{op} = \|A' \|_{op}$.
\end{lemma}

\begin{proof}
	There exists some $knl \times knl$ unitary permutation
matrix $P$,
\begin{equation}
  P (S_i \otimes \tilde{M} \otimes T_j) P^{\dag} = \tilde{M} \otimes S_i
  \otimes T_j .
\end{equation}
Obviously, there is $\| \tilde{M} \|_{op} = 1$, and
\begin{eqnarray}
  \|A' \|_{op} & = & \left\| \sum_{i, j} S_i \otimes \tilde{M} \otimes T_j
  \right\|_{op} \nonumber\\
  & = & \left\| P \left( \sum_{i, j} S_i \otimes \tilde{M} \otimes T_j
  \right) P^{\dag} \right\|_{op} \nonumber\\
  & = & \left\| \tilde{M} \otimes \sum_{i, j} S_i \otimes T_j \right\|_{op}
  \nonumber\\
  & = & \| \tilde{M} \otimes A\|_{op} = \| \tilde{M} \|_{op} \cdot \|A\|_{op}
  = \|A\|_{op} . 
\end{eqnarray}
\end{proof}

\begin{theorem}
  Consider a spectral triple $(\mathbb{M}_k (\mathbb{C}), \mathcal{H},
  \mathcal{D})$ with $\mathcal{H}=\mathbb{C}^{m k}$ and $\mathcal{D}= \sum_{i,
  j} S_i \otimes T_j$, $S_i$ are some $m' \times m'$ matrices, and $T_j$ are
  some $k' \times k'$ matrices with $k' \geqslant k$, and $mk = m' k'$. An
  element $a \in \mathbb{M}_k (\mathbb{C})$ acts on $\mathcal{H}$ through the
  diagonal representation $\pi (a) =\mathbb{I}_m \otimes a$. Then the
  corresponding spectral triple $(\mathbb{M}_k (\mathbb{C}), \mathcal{H}',
  \mathcal{D}')$ with
  \begin{equation}
    \mathcal{H}' =\mathbb{C}^n \otimes \mathcal{H}, \qquad \mathcal{D}' =
    \sum_{i, j} S_i \otimes \tilde{M} \otimes T_j
  \end{equation}
  has the same metric as $(\mathbb{M}_k (\mathbb{C}), \mathcal{H},
  \mathcal{D})$, where $M$ is any non-zero $n \times n$ Hermitian matrix,
  $\tilde{M} = \frac{M}{\|M\|_{op}}$, and an element $a \in \mathbb{M}_k
  (\mathbb{C})$ acts on $\mathcal{H}'$ through the diagonal representation
  $\pi' (a) =\mathbb{I}_{m n} \otimes a$.
\end{theorem}

\begin{proof}
	For simplicity, we only consider the case $k' = k$,
$m = m'$. There is
\begin{eqnarray}
  \|[\mathcal{D}', \pi' (e)]\|_{op} & = & \left\| \left[ \sum_{i, j} S_i
  \otimes \tilde{M} \otimes T_j, \mathbb{I}_{mn} \otimes e \right]
  \right\|_{op} \nonumber\\
  & = & \left\| \sum_{i, j} S_i \otimes \tilde{M} \otimes [T_j, e]
  \right\|_{op} \nonumber\\
  & = & \left\| \sum_{i, j} S_i \otimes [T_j, e] \right\|_{op}  \nonumber\\
  & = & \left\| \left[ \sum_{i, j} S_i \otimes T_j, \mathbb{I}_m \otimes e
  \right] \right\|_{op} \nonumber\\
  & = & \|[\mathcal{D}, \pi (e)]\|_{op} . 
\end{eqnarray}
So the spectral triples $(\mathbb{M}_k (\mathbb{C}), \mathcal{H},
\mathcal{D})$ and $(\mathbb{M}_k (\mathbb{C}), \mathcal{H}', \mathcal{D}')$
have the same ball condition,
\begin{equation}
  B' \equiv \{e \in \mathcal{A}: \|[\mathcal{D}', \pi' (e)]\|_{op} \leqslant 1\} =
  B \equiv \{e \in \mathcal{A}: \|[\mathcal{D}, \pi (e)]\|_{op} \leqslant 1\} .
\end{equation}
Therefore, the Connes spectral distances between the states $\rho_1, \rho_2
\in \mathbb{M}_k (\mathbb{C})$ in the spectral triple $(\mathbb{M}_k
(\mathbb{C}), \mathcal{H}', \mathcal{D}')$ are all equal to those of the
corresponding states in $(\mathbb{M}_k (\mathbb{C}), \mathcal{H},
\mathcal{D})$,
\begin{equation}
  d_{\mathcal{D}} (\rho_1, \rho_2) = \sup_{e \in B} | \mathrm{tr}_F (\Delta
  \rho \, e) | = \sup_{e \in B'} | \mathrm{tr}_F (\Delta \rho
  \, e) | = d_{\mathcal{D}'} (\rho_1, \rho_2) .
\end{equation}
\end{proof}

For example, it is easy to check that, in the following spectral triple, the
spectral distances between one-qubit states are equal to their quantum trace distances and also the Euclidean
distances of the corresponding Bloch vectors,
\begin{equation}
  \mathcal{A}=\mathbb{M}_2 (\mathbb{C}), \quad \mathcal{H}_8 =\mathbb{C}^8,
  \quad \mathcal{D}_8 = \frac{1}{4}  (\sigma_2 \otimes \tilde{M} \otimes
  \sigma_1 + \sigma_3 \otimes \tilde{M} \otimes \sigma_2 + \sigma_1 \otimes
  \tilde{M} \otimes \sigma_3),
\end{equation}
where $M$ is a non-zero $2\times 2$ Hermitian matrix, and
$\tilde{M} = \frac{M}{\|M\|_{op}}$.

The above concrete models can help us construct some other types of spectral triples with unitary invariant metric. These results can also help us understand the relations between the structure of Dirac operators and the properties of Connes spectral distance.

Furthermore, one can also use similar method to construct some spectral triples corresponding to higher-dimensional quantum states. But usually the calculations will be more cumbersome.



\section{Discussions and conclusions}\label{sec7}

In this paper, we study some properties of unitary invariance of Connes spectral distances between quantum states. We mainly focus on spectral triples with matrix algebras acting on finite dimensional Hilbert spaces via some linear representations. Furthermore, we also assume that the spectral distances are all finite. We derive some elementary properties of the Connes spectral distances and the optimal elements. The Lipschitz seminorms of the optimal elements are equal to $1$.
The unitary invariance of the spectral distances is equivalent to the unitary invariance of the Lipschitz seminorm.
If we assume the representation to be unital, then the spectral distances is also unitary invariant under some unitary transformations of the Dirac operators.

We also prove that there are some finite spectral triples in which the Lipschitz seminorms are equal to the operator norms.
In this case, the Connes spectral distances are just the normal quantum trace distances between quantum states.
If the representation is unital, then the only possible choice of the matrix algebra $\mathcal{A}$ is $\mathbb{M}_2(\mathbb{C})$. This just corresponds to the case of one-qubit states.
We also explicitly construct some concrete spectral triples for one-qubit states in which the Connes spectral distances are exactly the quantum trace distances.

These results are helpful for us to better study the relationships among Connes spectral distance, quantum trace distance and other quantum distance measures.
These concrete examples are significant for studies of geometric structures of finite spectral triples and mathematical relations of qubits and other quantum states in the framework of noncommutative geometry. These results can also be useful to
the studies of Connes spectral distances in other types of spectral triples.






\section*{Acknowledgments}
This work is partly supported by the
Guangdong Basic and Applied Basic Research Foundation (Grant No.
2024A1515010380), the 2024 Guangdong Province Education Science Planning
Project (Higher Education Special) (No. 2024GXJK455), the 2024 Guangdong
Higher Education Teaching Reform Project.



\begin{thebibliography}{100}

	\bibitem{Connes}A. Connes, \textit{Noncommutative geometry}
	(Academic Press, New York, 1994).
	
	\bibitem{Connes1}A. Connes, ``Compact metric spaces, Fredholm modules and
	hyperfiniteness.'' \textit{Ergodic Theory Dynam. Systems}
	\textbf{9} 207-220 (1989).
	
	\bibitem{Bimonte}G. Bimonte, F. Lizzi, G. Sparano, ``Distances on a
	lattice from non-commutative geometry.'' \textit{Phys. Lett. B}
	\textbf{341} 139-146 (1994).
	
	\bibitem{Dai}J. Dai, X.-C. Song, ``Connes' distance of one-dimensional
	lattices: general cases.'' \textit{Commun. Theor. Phys.}
	\textbf{36} 519 (2001).
	
	\bibitem{Cagnache}E. Cagnache, F. D'Andrea, P. Martinetti, J.-C. Wallet,
	``The spectral distance on the Moyal plane.'' \textit{J. Geom.
				Phys.} \textbf{61} 1881-1897 (2011).
	
	\bibitem{Wallet}J. C. Wallet, ``Connes distance by examples: Homothetic
	spectral metric spaces.'' \textit{Rev. Math. Phys.}
	\textbf{24} 1250027 (2012).
	
	\bibitem{Martinetti1}P. Martinetti, F. Mercati, L. Tomassini, ``Minimal
	length in quantum space and integrations of the line element in
	noncommutative geometry.'' \textit{Rev. Math. Phys.}
	\textbf{24} 1250010 (2012).
	
	\bibitem{Martinetti}P. Martinetti, L. Tomassini, ``Noncommutative geometry
	of the Moyal plane: translation isometries, Connes' distance on coherent
	states, Pythagoras equality.'' \textit{Commun. Math. Phys.}
	\textbf{323} 107--141 (2013).
	
	\bibitem{DAndrea}F. D'Andrea, F. Lizzi, J.C. Varilly, ``Metric properties
	of the fuzzy sphere.'' \textit{Lett. Math. Phys.}
	\textbf{103} 183--205 (2013).
	
	\bibitem{Pythagoras}F. D'Andrea, P. Martinetti, ``On Pythagoras theorem
	for products of spectral triples.'' \textit{Lett. Math. Phys.}
	\textbf{103} 469--492 (2013).
	
	\bibitem{Pythagoras1}F. D'Andrea, ``Pythagoras theorem in noncommutative
	geometry'' \textit{Contemporary Mathematics}
	\textbf{676} 175-210 (2016).
	
	\bibitem{Franco}N. Franco, J. C. Wallet, ``Metrics and causality on Moyal
	planes.'' \textit{Contemporary Mathematics}
	\textbf{676} 147-173 (2016).
	
	\bibitem{Scholtz}F. G. Scholtz, B. Chakraborty, ``Spectral triplets,
	statistical mechanics and emergent geometry in non-commutative quantum
	mechanics.'' \textit{J. Phys. A: Math. Theor.}
	\textbf{46} 085204 (2013).
	
	\bibitem{Chaoba}Y. Chaoba Devi, S. Prajapat, A. K. Mukhopadhyay, B.
	Chakraborty, F. G. Scholtz, ``Connes distance function on fuzzy sphere and
	the connection between geometry and statistics.'' \textit{J. Math.
				Phys.} \textbf{56} 041707 (2015).
	
	\bibitem{Revisiting}Y. Chaoba Devi, K. Kumar, B. Chakraborty, F. G.
	Scholtz, ``Revisiting Connes' finite spectral distance on noncommutative
	spaces: Moyal plane and fuzzy sphere.'' \textit{Int. J. Geo. Methods
				Mod. Phys.} \textbf{15} 1850204 (2018).
	
	\bibitem{Kumar}K. Kumar, B. Chakraborty, ``Spectral distances on the
	doubled Moyal plane using Dirac eigenspinors.'' \textit{Phys. Rev.
				D} \textbf{97} 086019 (2018).
	
	\bibitem{Barrett}J. W Barrett, P. Druce, L. Glaser, ``Spectral estimators
	for finite non-commutative geometries.'' \textit{J. Phys. A: Math.
				Theor.} \textbf{52} 275203 (2019).
	
	\bibitem{Chakraborty}A. Chakraborty, B. Chakraborty, ``Spectral distance
	on Lorentzian Moyal plane.'' \textit{Int. J. Geo. Methods Mod.
				Phys.} \textbf{17} 2050089 (2020).
				
	\bibitem{Lin1}B. S. Lin, T. H. Heng, ``Connes distance of $2 D$ harmonic
oscillators in quantum phase space.'' \textit{Chin. Phys. B}
\textbf{30} 110203 (2021).
	
	\bibitem{Lin2}B. S. Lin, T. H. Heng, ``Connes spectral distance and
	nonlocality of generalized noncommutative phase spaces.''
	\textit{Eur. Phys. J. Plus} \textbf{137} 899 (2022).
	
	\bibitem{Clare}P. Clare, C.-K. Li, E. Poon, E. Swartz, ``Noncommutative
	distances on graphs: An explicit approach via Birkhoff-James
	orthogonality.'' \textit{J. Geom. Phys.} \textbf{213}
	105483 (2025).

	\bibitem{Martinettip} P. Martinetti, ``From Monge to Higgs: a survey of distance computations in noncommutative geometry.'' \textit{Contemporary Mathematics.} \textbf{676} 1--46 (2016).


	\bibitem{Nielsen}M. A. Nielsen, I. L. Chuang, \textit{Quantum
				Computation and Quantum Information} (Cambridge University Press,
	Cambridge, 2000).
	
	\bibitem{Iochum}B. Iochum, T. Krajewski, P. Martinetti, ``Distances in
	finite spaces from noncommutative geometry.'' \textit{J. Geom.
				Phys.} \textbf{37} 100--125 (2001).
				
	\bibitem{Martinettif} P. Martinetti, ``Distances en geometrie non commutative.'' [arXiv:math-ph/0112038] (2001).
	

	
	\bibitem{Bassi} J. Bassi, R. Conti, ``On isometries of spectral triples associated to AF-algebras
	and crossed products.'' \textit{J. Noncommut. Geom.} \textbf{18} 547--566 (2024).

	\bibitem{Kumar1} K. Kumar, ``Calibrated Helstrom geometry on the Bloch ball via Connes spectral distance.'' [arXiv:2606.13824] (2026).

   \bibitem{LX} B. S. Lin, Z. H. Xu, J. H. Wang, H. L. Chen,
  ``Connes spectral distances, quantum discord and coherence of qubits.''
  \textit{Phys. Scr.} \textbf{101} 095204 (2026).
  
\end{thebibliography}

\end{document}
